\documentclass[10pt,a4paper]{amsart}
\usepackage[margin=19mm]{geometry}
\usepackage{amsmath,amssymb,mathtools}
\usepackage[scr=rsfs,cal=euler]{mathalfa}
\usepackage{xfrac}
\usepackage[unicode,hidelinks,bookmarks=false]{hyperref}
\newcommand{\ie}{i.e.\ }
\newcommand{\cf}{cf.\ }
\newcommand{\resp}{respectively\ }
\newcommand{\Subsection}{\subsection}
\providecommand{\nobreakdash}{\nobreak\hbox{-}}
\setlength{\emergencystretch}{2em}
\multlinegap0pt
\usepackage{amssymb}
\usepackage{cancel}
\usepackage[shortlabels]{enumitem}
\newtheorem{theorem}{Theorem}
\newtheorem{lemma}{Lemma}
\title{Explicit Formula Of The Critical Mass And The Energy Ground State Solution For The Mixed Local-Nonlocal Schrodinger Equation For The In-Between Critical Exponents Case}
\author{Yu Su, Hichem Hajaiej}
\date{Source: arXiv:2603.01138v1 (CC BY 4.0)}
\begin{document}
\maketitle

\noindent\emph{Source and licence.} Explicit Formula Of The Critical Mass And The Energy Ground State Solution For The Mixed Local-Nonlocal Schrodinger Equation For The In-Between Critical Exponents Case. Version arXiv:2603.01138v1 (CC BY 4.0), distributed by arXiv under the Creative Commons Attribution 4.0 International licence, http://creativecommons.org/licenses/by/4.0/, as recorded in the arXiv licence metadata for this exact version and shown on the abstract page https://arxiv.org/abs/2603.01138v1. Full licence notice: \texttt{licenses/arxiv-2603.01138-CC-BY-4.0.txt}.\par\smallskip

\noindent\textit{Editorial scope.} Counted unit: the article's lemma Lemma5.3 (source0.tex lines 1779--1808). The article's complete original proof of it is delivered with it (source0.tex lines 1810--1984, one \texttt{proof} environment of 175 lines). No mathematics is written, repaired or re-derived: the statement and the proof are the article's own lines, in the article's own notation, and no retained line is substituted or re-flowed.  Retained uncounted as the source-local context the proof needs: lines 130--136; lines 281--284; lines 1235--1243; lines 1294--1318; lines 1653--1669; lines 1709--1742.  Nothing was omitted from any retained span, so the package carries no omission record.  No retained line contains a citation, so the wrapper carries no bibliography and no bibliography entry.  No retained line contains a figure environment, an image-inclusion command or drawing code, so no image is delivered and the package declares no asset dependency.  Editorial formatting only, in addition: an AMS \texttt{amsart} preamble that restates the article's own declarations and macros used by the retained lines, namely the environments \texttt{theorem}, \texttt{lemma} on separate counters, together with the standard packages those lines need.  The retained environments print in this document as Theorem 1, Lemma 1 in order of appearance, whereas the article prints the same items with the numbering of its own full document.  The complete native source of the article as distributed in the arXiv e-print for this exact version is bundled unchanged as \texttt{sources/195544/source0.tex}.
\begin{equation}\label{MLN}
\begin{aligned}
-\Delta u
+(-\Delta)^{s}u+\omega u=|u|^{p-2}u, \ \ x\in \mathbb{R}^{N}.
\end{aligned}
\tag{MLN}
\end{equation}

\begin{theorem}\label{Theorem1.1}
Let $N\geqslant3$, $s\in(0,1)$ and  $p\in(2+\frac{4s}{N},2+\frac{4}{N})$.
Then there exists $c_{0}>0$ such that equation \eqref{MLN} has an energy ground state solution if and only if $c\geqslant c_{0}$.
\end{theorem}

\begin{equation}\label{4.1}
\begin{aligned}
\|u\|_{L^{p}(\mathbb{R}^{N})}^{p}
\leqslant C_{1,s,p}
\|u\|_{D^{s,2}(\mathbb{R}^{N})}^{\frac{2N+4-pN}{2(1-s)}}
\|u\|_{D^{1,2}(\mathbb{R}^{N})}^{\frac{pN-2N-4s}{2(1-s)}}
\|u\|_{L^{2}(\mathbb{R}^{N})}^{p-2}.
\end{aligned}
\end{equation}

\begin{lemma}\label{Lemma4.2}
Let $N\geqslant 3$, $s\in(0,1)$ and $p \in(2+\frac{4s}{N},2+\frac{4}{N})$.
Then there exists an optimizer $Q\in H^{1}(\mathbb{R}^{N})$ of inequality \eqref{4.1},
that is a weak solution of
\begin{equation}\label{4.4}
\begin{aligned}
-\Delta Q
+(-\Delta)^{s}Q+\omega Q=|Q|^{p-2}Q, \ \ x\in \mathbb{R}^{N},
\end{aligned}
\end{equation}
where $\omega>0,$ and
\begin{equation*}
\begin{aligned}
C_{1,s,p}
=
\frac
{\|Q\|_{L^{p}(\mathbb{R}^{N})}^{p}}
{
\|Q\|_{D^{s,2}(\mathbb{R}^{N})}^{\frac{2N+4-pN}{2(1-s)}}
\|Q\|_{D^{1,2}(\mathbb{R}^{N})}^{\frac{pN-2N-4s}{2(1-s)}}
\|Q\|_{L^{2}(\mathbb{R}^{N})}^{p-2}
}.
\end{aligned}
\end{equation*}
\end{lemma}

\begin{lemma}\label{Lemma5.1}
Let $N\geqslant 3$, $s\in(0,1)$ and $p \in(2+\frac{4s}{N},2+\frac{4}{N})$.
Set
\begin{equation*}
\begin{aligned}
\beta(Q)=\frac{\int_{\mathbb{R}^{N}}\int_{\mathbb{R}^{N}}\frac{|Q(x)-Q(y)|^{2}}{|x-y|^{N+2s}}\mathrm{d}x\mathrm{d}y}
{\int_{\mathbb{R}^{N}}|\nabla Q|^{2}\mathrm{d}x},
\end{aligned}
\end{equation*}
where $Q$ is defined in Lemma \ref{Lemma4.2}.
Then
\begin{equation*}
\begin{aligned}
\beta(Q)=\frac{2N+4-pN}{pN-2N-4s}.
\end{aligned}
\end{equation*}
\end{lemma}

\begin{lemma}\label{Lemma5.2}
Let $N\geqslant 3$, $s\in(0,1)$ and $p \in(2+\frac{4s}{N},2+\frac{4}{N})$.
Set
\begin{equation*}
\begin{aligned}
\gamma(Q)
=&\frac{\int_{\mathbb{R}^{N}}\int_{\mathbb{R}^{N}}\frac{|Q(x)-Q(y)|^{2}}{|x-y|^{N+2s}}\mathrm{d}x\mathrm{d}y}
{\int_{\mathbb{R}^{N}}|Q|^{p}\mathrm{d}x},
\end{aligned}
\end{equation*}
where $Q$ is defined in Lemma \ref{Lemma4.2}.
Then
\begin{equation*}
\begin{aligned}
\gamma(Q)=\frac{2N+4-pN}{2p(1-s)}.
\end{aligned}
\end{equation*}
And,
for $s\in(0,1)$, we have
\begin{equation*}
\begin{aligned}
\|Q\|^{2}_{D^{1,2}(\mathbb{R}^{N})}
=\frac{pN-2N-4s}{2p(1-s)}\int_{\mathbb{R}^{N}}|Q|^{p}\mathrm{d}x,
\end{aligned}
\end{equation*}
and
\begin{equation*}
\begin{aligned}
\|Q\|^{2}_{D^{s,2}(\mathbb{R}^{N})}
=\frac{2N+4-pN}{2p(1-s)}
\int_{\mathbb{R}^{N}}|Q|^{p}\mathrm{d}x.
\end{aligned}
\end{equation*}
\end{lemma}

\begin{lemma}\label{Lemma5.3}
Let $N\geqslant 3$, $s\in(0,1)$ and $p \in(2+\frac{4s}{N},2+\frac{4}{N})$.
Then
\begin{equation*}
\begin{aligned}
\frac{1}{2}
\int_{\mathbb{R}^{N}}
|\nabla Q|^{2}
\mathrm{d}x
+\frac{1}{2}
\int_{\mathbb{R}^{N}}
\int_{\mathbb{R}^{N}}
\frac{|Q(x)-Q(y)|^{2}}{|x-y|^{N+2s}}
\mathrm{d}x
\mathrm{d}y
-\frac{1}{p}
\int_{\mathbb{R}^{N}}
|Q|^{p}
\mathrm{d}x=0,
\end{aligned}
\end{equation*}
and
\begin{equation*}
\begin{aligned}
\int_{\mathbb{R}^{N}}
|Q|^{2}
\mathrm{d}x\geqslant c_{0}.
\end{aligned}
\end{equation*}
\end{lemma}

\begin{proof}
We know that
$Q$ satisifies the following Nehari identity
\begin{equation*}
\begin{aligned}
\int_{\mathbb{R}^{N}}
|\nabla Q|^{2}
\mathrm{d}x
+
\int_{\mathbb{R}^{N}}
\int_{\mathbb{R}^{N}}
\frac{|Q(x)-Q(y)|^{2}}{|x-y|^{N+2s}}
\mathrm{d}x
\mathrm{d}y
+
\omega
\int_{\mathbb{R}^{N}}
|Q|^{2}
\mathrm{d}x
-
\int_{\mathbb{R}^{N}}
|Q|^{p}
\mathrm{d}x
=0,
\end{aligned}
\end{equation*}
and the Pohoz\'{a}ev identity
\begin{equation*}
\begin{aligned}
&\frac{N-2}{2}
\int_{\mathbb{R}^{N}}
|\nabla Q|^{2}
\mathrm{d}x
+\frac{N-2s}{2}
\int_{\mathbb{R}^{N}}
\int_{\mathbb{R}^{N}}
\frac{|Q(x)-Q(y)|^{2}}{|x-y|^{N+2s}}
\mathrm{d}x
\mathrm{d}y\\
&+
\frac{N}{2}
\omega
\int_{\mathbb{R}^{N}}
|Q|^{2}
\mathrm{d}x
-\frac{N}{p}
\int_{\mathbb{R}^{N}}
|Q|^{p}
\mathrm{d}x
=0.
\end{aligned}
\end{equation*}
Then
\begin{equation*}
\begin{aligned}
\int_{\mathbb{R}^{N}}
|\nabla Q|^{2}
\mathrm{d}x
+s
\int_{\mathbb{R}^{N}}
\int_{\mathbb{R}^{N}}
\frac{|Q(x)-Q(y)|^{2}}{|x-y|^{N+2s}}
\mathrm{d}x
\mathrm{d}y
-\frac{(p-2)N}{2p}
\int_{\mathbb{R}^{N}}
|Q|^{p}
\mathrm{d}x=0
\end{aligned}
\end{equation*}
From Lemmas \ref{Lemma5.1} and \ref{Lemma5.2},
we get that
\begin{equation*}
\begin{aligned}
&\frac{1}{2}
\int_{\mathbb{R}^{N}}
|\nabla Q|^{2}
\mathrm{d}x
+(s-\frac{1}{2})
\int_{\mathbb{R}^{N}}
\int_{\mathbb{R}^{N}}
\frac{|Q(x)-Q(y)|^{2}}{|x-y|^{N+2s}}
\mathrm{d}x
\mathrm{d}y
-\frac{1}{p}
(\frac{(p-2)N}{2}-1)
\int_{\mathbb{R}^{N}}
|Q|^{p}
\mathrm{d}x\\
=&
\left[\frac{1}{2}\frac{pN-2N-4s}{2N+4-pN}
+s-\frac{1}{2}
-\frac{1}{p}
(\frac{(p-2)N}{2}-1)
\frac{2p(1-s)}{2N+4-pN}\right]
\int_{\mathbb{R}^{N}}
\int_{\mathbb{R}^{N}}
\frac{|Q(x)-Q(y)|^{2}}{|x-y|^{N+2s}}
\mathrm{d}x
\mathrm{d}y\\
=&0.
\end{aligned}
\end{equation*}
Therefore,
\begin{equation*}
\begin{aligned}
0=&
\int_{\mathbb{R}^{N}}
|\nabla Q|^{2}
\mathrm{d}x
+s
\int_{\mathbb{R}^{N}}
\int_{\mathbb{R}^{N}}
\frac{|Q(x)-Q(y)|^{2}}{|x-y|^{N+2s}}
\mathrm{d}x
\mathrm{d}y
-\frac{(p-2)N}{2p}
\int_{\mathbb{R}^{N}}
|Q|^{p}
\mathrm{d}x\\
=&\frac{1}{2}
\int_{\mathbb{R}^{N}}
|\nabla Q|^{2}
\mathrm{d}x
+(s-\frac{1}{2})
\int_{\mathbb{R}^{N}}
\int_{\mathbb{R}^{N}}
\frac{|Q(x)-Q(y)|^{2}}{|x-y|^{N+2s}}
\mathrm{d}x
\mathrm{d}y
-\frac{1}{p}
(\frac{(p-2)N}{2}-1)
\int_{\mathbb{R}^{N}}
|Q|^{p}
\mathrm{d}x\\
&+\frac{1}{2}
\int_{\mathbb{R}^{N}}
|\nabla Q|^{2}
\mathrm{d}x
+\frac{1}{2}
\int_{\mathbb{R}^{N}}
\int_{\mathbb{R}^{N}}
\frac{|Q(x)-Q(y)|^{2}}{|x-y|^{N+2s}}
\mathrm{d}x
\mathrm{d}y
-\frac{1}{p}
\int_{\mathbb{R}^{N}}
|Q|^{p}
\mathrm{d}x\\
=&\frac{1}{2}
\int_{\mathbb{R}^{N}}
|\nabla Q|^{2}
\mathrm{d}x
+\frac{1}{2}
\int_{\mathbb{R}^{N}}
\int_{\mathbb{R}^{N}}
\frac{|Q(x)-Q(y)|^{2}}{|x-y|^{N+2s}}
\mathrm{d}x
\mathrm{d}y
-\frac{1}{p}
\int_{\mathbb{R}^{N}}
|Q|^{p}
\mathrm{d}x.
\end{aligned}
\end{equation*}
From Theorem \ref{Theorem1.1},
we know that
\begin{equation*}
\begin{aligned}
\int_{\mathbb{R}^{N}}
|Q|^{2}
\mathrm{d}x\geqslant c_{0}.
\end{aligned}
\end{equation*}
\end{proof}

\end{document}
